% questions.tex —— 1955 年·全国卷 · 题目内容片段（无答案）
% 由 tools/generate_projects.py 依来源总账自 1955/1955_数学真题_试题.md 生成，勿手改（改源 MD 后重跑）


\par\medskip

\gkpart{一、小题}

\begin{question}
\gkti{1}（甲）以二次方程 $x^2-3x-1=0$ 的两根的平方为两根作一二次方程．

\end{question}
\begin{question}
\gkti{2}（乙）等腰三角形一腰的长是底边的 $4$ 倍，求这三角形各角的余弦．

\end{question}
\begin{question}
\gkti{3}（丙）已知正四棱锥底边的长为 $a$，侧棱与底面的交角为 $45^\circ$，求这棱锥的高．

\end{question}
\begin{question}
\gkti{4}（丁）写出：二面角的平面角的定义．

\par\medskip

\end{question}
\gkpart{二、大题}

\begin{question}
\gkti{5}（二）求 $b$，$c$，$d$ 的值，使多项式 $x^3+bx^2+cx+d$ 适合下列三条件：（1）被 $x-1$ 整除；（2）被 $x-3$ 除时余 $2$；（3）被 $x+2$ 除与被 $x-2$ 除时余数相等．

\end{question}
\begin{question}
\gkti{6}（三）由直角三角形勾上一点 $D$ 作弦 $AB$ 的垂线交弦于 $E$、股的延长线于 $F$、外接圆周于 $Q$，求证：$EQ$ 为 $EA$ 与 $EB$ 的比例中项又为 $ED$ 与 $EF$ 的比例中项．

\end{question}
\begin{question}
\gkti{7}（四）解方程 $\cos2x=\cos x+\sin x$，求 $x$ 的通值．

\end{question}
\begin{question}
\gkti{8}（五）一三角形三边的长成等差级数，其周长为 $12$ 尺，面积为 $6$ 平方尺，求证这三角形为一直角三角形．

\end{question}